In this paper, we demonstrate that structured prompt engineering exerts a
decisive impact on the capability of Large Language Models to simulate empirical
public opinion distributions in Brazil. Our findings establish that cumulative
prompt refinement progressively minimizes distributional divergence, achieving
optimal alignment ($\text{JSD} = 0.121$) through the synergy of demographic
feature selection, structural role separation, Chain-of-Thought deliberation,
and in-context calibration. Furthermore, our innovative benchmarking between
culturally specialized Brazilian architectures (Sabiá-4 and Sabiazinho-4) and
general-purpose global baselines (GPT-5-mini and LLaMA~3.2~3B) reveals that
domain-tailored models exhibit superior sociocultural fidelity by accurately
capturing empirical civic skepticism, whereas global baselines succumb to severe
social desirability and institutional optimism biases. Within this framework,
Explainable AI (XAI) auditing proves indispensable for detecting cognitive
dissonance and verifying persona behavioral consistency.

The authors acknowledge, however, certain methodological limitations of these
findings. Most notably, the empirical benchmark relies on a single
cross-sectional survey instrument (CESOP/IPEC) centered on democratic
perceptions. While this instrument provides rich structural variance across
multi-alternative policy, institutional, and engagement items, fully
consolidating the external validity of silicon sampling warrants replicating
this evaluative pipeline across broader thematic domains and longitudinal
polling cycles.

Promising avenues for future research involve deploying Retrieval-Augmented
Generation (RAG) to dynamically inject real-time regional context and exploring
supervised fine-tuning of open-weight architectures on authentic citizen
rationales. Extending this probabilistic framework across additional open-weight
foundations and sampling temperatures will further strengthen simulation
robustness. Ultimately, this work demonstrates the strong potential of Large
Language Models for public opinion simulation and establishes solid
methodological directions for their continued development.

